% chapters/02/section-02.tex
\section{Model Waktu: Bingkai dan Langkah Simulasi}
\label{sec:model-waktu}

Bagian ini membahas persoalan yang paling sering merugikan pemula dan yang
paling jarang disadari: pengelolaan waktu. Hampir semua keluhan bahwa sebuah
permainan ``terasa aneh'' pada satu komputer tetapi tidak pada komputer lain
berpangkal pada bagian ini.

\subsection{Mengapa Waktu Menjadi Persoalan}

Komputer menjalankan perintah jauh lebih cepat daripada yang dapat diamati
manusia. Sebuah komputer dapat menyelesaikan puluhan juta operasi dalam satu
detik, sedangkan mata manusia mulai melihat urutan gambar sebagai gerak yang
mulus pada sekitar dua puluh sampai enam puluh gambar per detik. Karena itu
program permainan tidak berjalan secepat mungkin; ia berjalan menurut irama
yang ditentukan.

Irama tersebut tidak dapat diandalkan untuk mengukur waktu. Program yang
memindahkan objek sejauh lima piksel pada setiap pengulangan akan
menghasilkan gerak yang berbeda-beda, bergantung pada berapa kali pengulangan
tersebut selesai dalam satu detik. Pada komputer yang cepat pengulangannya
lebih sering, sehingga objek bergerak lebih jauh. Akibatnya, permainan menjadi
lebih sulit pada komputer yang lebih cepat, dan kecepatan permainan menjadi
tidak dapat diramalkan.

\subsection{Bingkai dan Laju Bingkai}

Satu putaran lengkap siklus permainan disebut satu bingkai (\istilah{frame}),
dan banyaknya bingkai yang diselesaikan dalam satu detik disebut laju bingkai
(\istilah{frame rate}), dinyatakan dalam satuan bingkai per detik
(\istilah{frames per second}, FPS). Laju bingkai bukan sekadar angka
kecepatan; ia adalah anggaran waktu: pada 60 FPS, seluruh pekerjaan satu
bingkai harus selesai dalam sekitar 16,7 milidetik.

Anggaran tersebut menuntut sikap yang berbeda dalam menulis program.
Menambah satu perhitungan berat pada tahap simulasi berarti mengurangi jatah
waktu tahap lain. Ketika anggaran terlampaui, laju bingkai turun, dan
penurunan itu langsung terasa oleh pemain. Karena itu program permainan lazim
membatasi laju bingkainya pada nilai tertentu, misalnya 60 FPS, agar
pemakaian daya tidak berlebihan tanpa manfaat yang sepadan
\cite{sweigart-mg,pygame-docs}.

\subsection{Selang Waktu sebagai Jembatan}

Penyelesaian persoalan di atas terletak pada satu gagasan sederhana: jangan
menyatakan gerak dalam satuan per bingkai, melainkan dalam satuan per detik.
Apabila kecepatan objek dinyatakan sebagai 200 piksel per detik, maka jarak
yang harus ditempuh pada satu bingkai bergantung pada lama bingkai itu
berlangsung. Besaran lama satu bingkai inilah yang disebut selang waktu
(\istilah{delta time}) dan lazim dilambangkan $\Delta t$.

Hubungan ketiga besaran tersebut dapat dituliskan sebagai

\begin{equation}
\label{eq:gerak-dt}
\text{perpindahan} = \text{kecepatan} \times \Delta t .
\end{equation}

Apabila satu bingkai berlangsung 16,7 milidetik, perpindahannya sekitar 3,3
piksel; apabila bingkai berikutnya berlangsung dua kali lebih lama, yaitu 33,3
milidetik, perpindahannya menjadi 6,7 piksel. Dengan cara ini, objek menempuh
jarak yang sama dalam satu detik, terlepas dari berapa banyak bingkai yang
dihasilkan komputer \cite{fiedler-timestep}.

\subsection{Langkah Tetap dan Langkah Berubah}

Sampai di sini, panjang langkah simulasi mengikuti lama bingkai yang
sebenarnya. Cara ini disebut langkah berubah (\istilah{variable timestep}) dan
paling mudah diterapkan. Kelemahannya, panjang langkah yang selalu berubah
membuat hasil simulasi sulit diulang secara persis dan membuat perhitungan
fisika kurang stabil ketika komputer sesekali melambat.

Alternatifnya adalah langkah tetap (\istilah{fixed timestep}). Simulasi selalu
maju dengan panjang langkah yang sama, misalnya 1/60 detik, sedangkan selisih
waktu nyata yang belum terpakai ditabung pada sebuah penampung
(\istilah{accumulator}). Apabila penampung telah mencapai satu panjang
langkah, simulasi dimajukan satu langkah dan penampung dikurangi. Jika
komputer melambat, simulasi dapat dimajukan beberapa langkah sekaligus;
apabila komputer terlalu cepat, simulasi menunggu sampai waktu nyatanya tiba.
Perbandingan kedua cara ini diperlihatkan pada Gambar~\ref{fig:waktu-bingkai}.

\begin{figure}[htbp]
\centering
% chapters/02/figures/waktu-bingkai.tikz
% Perbandingan langkah berubah (variable timestep) dan langkah tetap
% (fixed timestep). Jarak antar-tanda tegak menyatakan lama satu bingkai.
\begin{tikzpicture}[font=\scriptsize]
  % ---------- Langkah berubah ----------
  \node[anchor=west, ubPrimary, font=\footnotesize\bfseries] at (0,3.45)
    {Langkah berubah: panjang langkah mengikuti lama bingkai};
  \draw[ubAccent, thick] (0,3.0) -- (11.2,3.0);
  \foreach \x in {0,0.9,2.6,3.3,5.0,5.6,7.3,8.0,8.7,10.4,11.2}
    \draw[ubAccent, thick] (\x,2.85) -- (\x,3.15);
  \node[anchor=west, ubRed] at (0.75,2.45) {$\Delta t$ besar};
  \node[anchor=west, ubGreen] at (6.55,2.45) {$\Delta t$ kecil};
  \draw[ubRed, -{Stealth[length=1.8mm]}] (1.45,2.62) -- (1.05,2.86);
  \draw[ubGreen, -{Stealth[length=1.8mm]}] (7.15,2.62) -- (7.62,2.86);
  \node[anchor=west, font=\scriptsize\itshape] at (0,1.85)
    {\istilah{Delta time} menjaga kecepatan objek tetap sama di setiap komputer.};

  % ---------- Langkah tetap ----------
  \node[anchor=west, ubPrimary, font=\footnotesize\bfseries] at (0,1.05)
    {Langkah tetap: panjang langkah selalu sama, sisa waktu ditabung};
  \draw[ubAccent, thick] (0,0.6) -- (11.2,0.6);
  \foreach \x in {0,1.12,2.24,3.36,4.48,5.60,6.72,7.84,8.96,10.08,11.2}
    \draw[ubAccent, thick] (\x,0.45) -- (\x,0.75);
  \node[anchor=west, font=\scriptsize\itshape] at (0,-0.15)
    {Selisih waktu yang belum terpakai disimpan pada penampung (\istilah{accumulator}).};

  % ---------- Sumbu waktu ----------
  \draw[ubLine, -{Stealth[length=2mm]}] (0,-0.85) -- (11.4,-0.85);
  \node[anchor=north, ubLine] at (5.6,-0.9) {waktu nyata};
\end{tikzpicture}

\caption{Perbandingan langkah berubah dan langkah tetap pada garis waktu}
\label{fig:waktu-bingkai}
\end{figure}

Pilihan di antara keduanya bergantung pada kebutuhan. Permainan yang
mengandalkan ketepatan fisika dan hasil yang dapat diulang, misalnya permainan
yang menyimpan dan memutar kembali rekaman gerak, lebih cocok memakai langkah
tetap. Permainan sederhana yang hanya memerlukan gerak halus cukup memakai
langkah berubah. Yang perlu Anda hindari bukan salah satu pilihan itu,
melainkan ketiadaan pilihan: menulis gerak tanpa memperhitungkan waktu sama
sekali.

\begin{catatabox}
Mesin permainan seperti Unity menawarkan kedua cara tersebut sekaligus.
Perhitungan yang berkaitan dengan tampilan dijalankan pada tahap yang
mengikuti lama bingkai, sedangkan perhitungan fisika dijalankan pada tahap
dengan langkah tetap. Kedua tahap itu akan Anda temui pada Bab~XIII sebagai
\code{Update} dan \code{FixedUpdate}.
\end{catatabox}

\subsection{Akibat pada Perilaku Permainan}

Keputusan tentang model waktu bukan urusan teknis semata; ia memengaruhi
pengalaman bermain. Tiga akibat berikut pantas diperhatikan sejak awal.

\begin{enumerate}
  \item \textbf{Keadilan permainan.} Apabila gerak bergantung pada laju
        bingkai, pemain dengan komputer lebih cepat memperoleh keuntungan yang
        tidak dirancang.
  \item \textbf{Kemudahan pengujian.} Langkah tetap membuat simulasi berjalan
        sama pada setiap percobaan. Kesamaan itu memudahkan Anda memastikan
        apakah sebuah perubahan kode memperbaiki atau justru merusak perilaku
        permainan.
  \item \textbf{Kestabilan perhitungan.} Langkah yang terlalu panjang membuat
        perhitungan fisika dapat melampaui batas dan menghasilkan perilaku yang
        tidak diinginkan, misalnya objek menembus dinding. Persoalan ini
        dibahas pada Bagian~\ref{sec:simulasi-diskret}.
\end{enumerate}
